\section{研究方法}
\label{sec:methodology}

本研究比較六種結合 Kalman filter、ARIMA 與機器學習之時間序列點預測及區間預測方法。M1--M5 依洪瑋翎之研究流程實作；M6 為本研究提出之 Kalman--ARIMA--ANN Out-of-Time Adaptive EnbPI-type 方法。除另行說明之控制實驗外，各方法皆保留原始模型結構與超參數。所有方法均採 sequential one-step-ahead evaluation：時間 $t$ 的預測完成後才揭露 $Y_t$，且 $Y_t$ 僅能用於時間 $t+1$ 以後的預測。

\subsection{M6：Kalman--ARIMA--ANN Out-of-Time Adaptive EnbPI}
\label{subsec:m6-method}

\subsubsection{訓練資料切割與整體估計流程}

首先先切割資料，在每一個時間 fold 內執行分解與模型訓練。完整資料分為 Training 與 Testing：
\[
    \underbrace{Y_1,\ldots,Y_T}_{\text{training}}
    \quad\Big|\quad
    \underbrace{Y_{T+1},\ldots,Y_{T+H}}_{\text{final test}}.
\]
在 Training period 內，將前 $\rho$ 筆資料定義為 initial training prefix，並將其後的時間點切成 $K$ 個彼此不重疊且連續的 future validation blocks：
\[
    I_1,I_2,\ldots,I_K.
\]
對於第 $k$ 個 future block 的第一個時間點。所有步驟，包括 Kalman 分解、ARIMA order selection、ANN iteration selection 與 $B$ 組 bootstrap models 的訓練，都只能使用該時間點以前的資料。
模型隨後對 $I_k$ 逐期進行 one-step-ahead prediction。完成整個 block 且其中真值均已揭露後，該 block 才併入下一個 fold 的 training prefix。資料流程可寫成
\[
\begin{aligned}
Y_{1:325}
&\longrightarrow \text{訓練模型並預測 }I_1,\\
Y_{1:\max(I_1)}
&\longrightarrow \text{重新分解、選模、訓練並預測 }I_2,\\
Y_{1:\max(I_2)}
&\longrightarrow \text{重新分解、選模、訓練並預測 }I_3,\\
Y_{1:\max(I_3)}
&\longrightarrow \text{重新分解、選模、訓練並預測 }I_4.
\end{aligned}
\]

\subsubsection{Kalman 分解}
在任一時間 $t$，僅使用時間 $t$ 之前的資料執行 Kalman filter，得到低頻估計$\widehat L_t^{KF}$，
並將高頻成分定義為觀測值扣除低頻估計：
\[\widehat H_t^{KF}=Y_t-\widehat L_t^{KF}\]
因此對每個時間點皆有
\[
    Y_t=\widehat L_t^{KF}+\widehat H_t^{KF}.
\]

\subsubsection{高頻 ANN supervised windows}
令 lag window 為 $w$。對每個 $t=w+1,\ldots,m$，以過去 $w$ 期的 Kalman low、Kalman high 與原始觀測建立 ANN 輸入：
\[
    X_t=(
    \widehat L_{t-w:t-1}^{KF},
    \widehat H_{t-w:t-1}^{KF},
    Y_{t-w:t-1}
    )
\]
ANN 的預測目標為當期 Kalman high：
\[
    Z_t=\widehat H_t^{KF}.
\]
因此一筆 supervised sample 為
\[
    D_t=(X_t,Z_t).
\]

\subsubsection{低頻 base ARIMA 與 order selection}

在每個 fold 中，只使用該 validation block 開始以前的 Kalman low series，於候選集合
\[
    \mathcal P=\{(p,0,q):0\leq p\leq4,\ 0\leq q\leq4\}
\]
中以 BIC 選擇一次共同 order：
\[
    (p^\ast,0,q^\ast)
    =\arg\min_{(p,0,q)\in\mathcal P}\operatorname{BIC}(p,0,q).
\]
使用此 order 在依時間排列的 Kalman low series 上擬合 base ARIMA，並取得 one-step fitted values $\widehat L_t^{\mathrm{base}}$ 及殘差
\[
    u_t=\widehat L_t^{KF}-\widehat L_t^{\mathrm{base}}.
\]
為避免將 base model 的有限樣本平均誤差重複加入 bootstrap series，先將殘差中心化：
\[
    \widetilde u_t=u_t-\overline u.
\]

\subsubsection{Moving-block bootstrap 與 $B$ 組 hybrid models}

令 supervised samples 數量為 $n$，bootstrap block length 為
\[
    \ell=\max\{2,\operatorname{round}(\sqrt n)\}.
\]
第 $b$ 次 bootstrap 隨機選取
\[
    K_b=\frac{n}{\ell}
\]
個 block 起點，每個起點向後取得 $\ell$ 個連續索引，串接後截取前 $n$ 個索引，形成
\[
    S_b=(s_{b1},\ldots,s_{bn}).
\]
此程序重複 $B$ 次。

\paragraph{第 $b$ 個 ARIMA。}

ARIMA 分支抽取 base ARIMA residual blocks，並將抽得的 residual 逐點加回 base fitted values：
\[
    L_t^{\ast(b)}
    =\widehat L_t^{\mathrm{base}}+\widetilde u_{s_{bt}}.
\]
第 $b$ 個 ARIMA 使用同一個已選 order $(p^\ast,0,q^\ast)$，並重新估計模型參數：
\[
    \mathcal A_b
    =\operatorname{ARIMA}\!\left(L^{\ast(b)};(p^\ast,0,q^\ast)\right).
\]

\paragraph{第 $b$ 個 ANN。}

ANN 分支使用相同索引 $S_b$ 抽取按時間排列的 supervised samples：
\[
    \mathcal D_b^\ast=\{(X_{s_{bj}},Z_{s_{bj}}):j=1,\ldots,n\}.
\]
每個 $D_t=(X_t,Z_t)$ 是不可拆分的一筆訓練樣本；moving-block bootstrap 抽取的是相鄰 target times 所對應的整筆 supervised samples。

輸入 $X$ 與 target $Z$ 均在第 $b$ 個 bootstrap training set 內標準化。ANN 採兩層 hidden layers $(32,16)$、ReLU activation 與 MSE loss：
\[
    \mathcal L_b(\theta)
    =\frac{1}{n}\sum_{j=1}^{n}
    \left[Z_{s_{bj}}-g_{\theta}(X_{s_{bj}})\right]^2.
\]
最大 iteration 由候選集合 $\{125,250,500\}$ 中，透過 bootstrap sample 內依原始時間排序的 rolling validation 選定。由於各 replicate 的樣本、重複次數、random seed 與可能選到的 iteration 不同，最終得到 $B$ 個 ANN：
\[
    \mathcal N_1,\ldots,\mathcal N_B.
\]

\paragraph{Hybrid aggregation。}

同一次 bootstrap 的 $\mathcal A_b$ 與 $\mathcal N_b$ 配成第 $b$ 個 hybrid predictor。使用 mean aggregation：
\[
    \widehat L_t^{\mathrm{raw}}
    =\frac{1}{B}\sum_{b=1}^{B}\widehat L_{t,b},
    \qquad
    \widehat H_t^{\mathrm{raw}}
    =\frac{1}{B}\sum_{b=1}^{B}\widehat H_{t,b},
\]
並得到未校正的合併點預測
\[
    \widetilde Y_t
    =\widehat L_t^{\mathrm{raw}}+\widehat H_t^{\mathrm{raw}}.
\]

\subsubsection{各 future blocks 的 out-of-time residuals}

依照前述 rolling-origin 切割，第 $k$ 個 fold 的 predictor 只使用該 fold 前的數值完成分解、model selection 與 fitting。為避免先使用完整 training series 選擇模型，再回頭校準同一段資料所造成的資訊洩漏，第 $k$ 個 future block 的每一筆 residual 都必須由此 predictor 產生。

接著對 $t\in I_k$ 進行 sequential one-step-ahead prediction：模型參數在該 block 內固定，但 $Y_t$ 只能在完成時間 $t$ 的預測後，加入歷史以供 $t+1$ 使用。令
\[
    \widehat Y_{t\mid t-1}
    =\widehat L_{t\mid t-1}+\widehat H_{t\mid t-1}
\]
表示只依賴時間 $t-1$ 以前資訊的合併點預測，並定義 calibration residual
\[
    r_t^{\mathrm{cal}}
    =Y_t-\widehat Y_{t\mid t-1}.
\]
完成第 $k$ 個 block 後，該 block 已揭露的真實資料才可加入 expanding training prefix，再重新執行整套 model selection 與 fitting 以預測下一個 block。

最終 training calibration sequence 為
\[
    \mathcal E_{\mathrm{cal}}
    =\{r_t^{\mathrm{cal}}:t\in I_1\cup\cdots\cup I_K\}.
\]
\subsubsection{Adaptive residual window selection}

為容許預測誤差分布隨時間改變，考慮候選 residual-window lengths
\[
    \mathcal W=\{100,200,300,400,\mathrm{all}\}.
\]
對 training cross-fit residual sequence 的後 $40\%$ 進行 prequential validation。在每個 validation residual 出現以前，只使用先前 residuals 建立區間，並以 Gneiting--Raftery interval score 評估：
\[
S_{\alpha}(l,u;y)
=(u-l)
+\frac{2}{\alpha}(l-y)\mathbf 1\{y<l\}
+\frac{2}{\alpha}(y-u)\mathbf 1\{y>u\}.
\]
選擇平均 interval score 最小的 window：
\[
    \widehat W
    =\arg\min_{W\in\mathcal W}
    \frac{1}{|\mathcal V|}
    \sum_{j\in\mathcal V}S_{\alpha}(l_{j,W},u_{j,W};e_j^{CF}).
\]
此選擇完全在 training cross-fit residuals 上完成，不使用 final test outcomes。進入測試階段後，$\widehat W$ 的數值固定，但 residual pool 的內容會隨新真值揭露逐步更新。

\subsubsection{Final sequential point and interval prediction}

完成 calibration 後，使用完整 training set $Y_{1:T}$ 重新執行 Kalman 分解、ARIMA order selection、ANN iteration selection，以及 $B$ 組 hybrid models 的訓練。初始測試 residual pool 為 training cross-fit residuals 的最後 $\widehat W$ 筆：
\[
    \mathcal E_T
    =\operatorname{tail}_{\widehat W}(\mathcal E^{CF}).
\]

對每一個測試時間 $t=T+1,\ldots,T+H$，先計算合併 raw forecast $\widetilde Y_t$。並使用 combined bias correction：
\[
    b_t=\frac{1}{|\mathcal E_{t-1}|}
    \sum_{e\in\mathcal E_{t-1}}e,
    \qquad
    \widehat Y_t=\widetilde Y_t+b_t.
\]
接著在尚未觀察 $Y_t$ 前，以目前 residual pool 選擇最短 residual interval 的 $\beta$：
\[
\widehat\beta_t
=\arg\min_{\beta\in[0,\alpha]}
\left\{
Q_{1-\alpha+\beta}(\mathcal E_{t-1}-b_t)
-Q_{\beta}(\mathcal E_{t-1}-b_t)
\right\}
\]
令
\[
q_t^-=Q_{\widehat\beta_t}(\mathcal E_{t-1}-b_t),
\qquad
q_t^+=Q_{1-\alpha+\widehat\beta_t}(\mathcal E_{t-1}-b_t),
\]
則輸出的 prediction interval 為
\[
    C_t^{1-\alpha}
    =[\widehat Y_t+q_t^-,\widehat Y_t+q_t^+].
\]
當 $\widehat Y_t$ 與 $C_t^{1-\alpha}$ 均輸出後，才揭露 $Y_t$ 並計算新的 raw residual：
\[
    e_t=Y_t-\widetilde Y_t.
\]
殘差池更新為
\[
    \mathcal E_t
    =\operatorname{tail}_{\widehat W}
    \left(\mathcal E_{t-1}\cup\{e_t\}\right).
\]
因此 $e_t$ 只能影響時間 $t+1$ 以後的區間，不會回頭改變時間 $t$ 已輸出的結果。

\begin{algorithm}[htbp]
\caption{Kalman--ARIMA--ANN Out-of-Time Adaptive EnbPI}
\label{alg:kf-oot-enbpi}

\begin{spacing}{1}

\begin{algorithmic}[1]

\Require
訓練序列 $\{Y_t\}_{t=1}^{T}$；
測試長度 $H$；
錯誤率 $\alpha$；
bootstrap 數量 $B$；
lag window $w$；
block length $\ell$；
future-block 數量 $K$；
residual-window 候選集合 $\mathcal W$

\Ensure
測試期點預測 $\widehat Y_{t\mid t-1}$
與預測區間 $C_t^{1-\alpha}$

\State
令 $a_0=\frac{T}{2}$；
以前 $a_0$ 筆為 initial training，
後半段切成連續 validation blocks $I_1,\ldots,I_K$

\State
初始化 calibration residual sequence $\mathcal E_{\mathrm{cal}}\gets\varnothing$

\For{$k=1,\ldots,K$}

    \State
    令 $s_k=\min(I_k)$，
    第 $k$ 個 fold 只使用 $Y_{1:s_k-1}$ 訓練

    \State
    Kalman 分解：
    $\widehat L_t^{KF}=KF(Y_t)$，
    $\widehat H_t^{KF}=Y_t-\widehat L_t^{KF}$

    \State
    建立 ANN supervised samples 
    $X_t=(\widehat L_{t-w:t-1}^{KF},
    \widehat H_{t-w:t-1}^{KF},Y_{t-w:t-1})$，
    target 為 $Z_t=\widehat H_t^{KF}$

    \State
    以 BIC 選擇一次 ARIMA order，
    擬合 base ARIMA 並取得 residuals $u_t$

    \For{$b=1,\ldots,B$}

        \State
        以 moving-block bootstrap 產生 $S_b$；
        用 $u_t[S_b]$ 擬合相同 order 的 ARIMA$_b$；
        用 $(X_t,Z_t)[S_b]$ 與 MSE 訓練 ANN$_b$

    \EndFor

    \For{$t\in I_k$}

        \State
        使用 $Y_{1:t-1}$ 計算
        $\widehat Y_{t\mid t-1}
        =B^{-1}\sum_{b=1}^{B}
        (\widehat L_{t,b}+\widehat H_{t,b})$

        \State
        揭露 $Y_t$，並令
        $\mathcal E_{\mathrm{cal}}
        \gets\mathcal E_{\mathrm{cal}}
        \cup\{Y_t-\widehat Y_{t\mid t-1}\}$

    \EndFor
\EndFor

\State
以 training-only interval score 選擇
$\widehat W
=\arg\min_{W\in\mathcal W}
\operatorname{Score}(W;\mathcal E_{\mathrm{cal}})$

\State
令 $\mathcal E
=\operatorname{tail}_{\widehat W}
(\mathcal E_{\mathrm{cal}})$，
並以完整 $Y_{1:T}$ 重新訓練最終 $B$ 組模型

\For{$t=T+1,\ldots,T+H$}

    \State
    計算
    $\widetilde Y_t
    =B^{-1}\sum_{b=1}^{B}
    (\widehat L_{t,b}+\widehat H_{t,b})$

    \State
    令 $b_t=\operatorname{mean}(\mathcal E)$，
    $\widehat Y_{t\mid t-1}=\widetilde Y_t+b_t$

    \State
    選擇
    $\widehat\beta_t
    =\arg\min_{\beta\in[0,\alpha]}
    \{Q_{1-\alpha+\beta}(\mathcal E-b_t)
    -Q_\beta(\mathcal E-b_t)\}$

    \State
    輸出
    $C_t^{1-\alpha}
    =[\widehat Y_{t\mid t-1}
    +Q_{\widehat\beta_t}(\mathcal E-b_t),
    \widehat Y_{t\mid t-1}
    +Q_{1-\alpha+\widehat\beta_t}(\mathcal E-b_t)]$

    \State
    揭露 $Y_t$，更新
    $\mathcal E\gets
    \operatorname{tail}_{\widehat W}
    (\mathcal E\cup\{Y_t-\widetilde Y_t\})$

\EndFor

\State \Return
$\{\widehat Y_{t\mid t-1},C_t^{1-\alpha}\}_{t=T+1}^{T+H}$

\end{algorithmic}
\end{spacing}
\end{algorithm}

\subsection{M1--M5 比較方法之共同流程}
\label{subsec:m1-m5-common}

M1--M5 先對完整觀測序列進行因果 Kalman filter 分解，令
\[
    Y_t=\widehat L_t^{KF}+\widehat H_t^{KF},
    \qquad
    \widehat H_t^{KF}=Y_t-\widehat L_t^{KF}.
\]
Kalman filter 的 process variance 與 measurement variance 分別設為 $0.5$ 與 $10$。低頻分支只使用 training period 的 $\widehat L_t^{KF}$，在
\[
    \{(p,0,q):p,q\in\{0,1,2,3,4\}\}
\]
中以 AIC 選擇 ARIMA order；進入測試期後採 rolling one-step-ahead prediction，時間 $t$ 的真值揭露後才加入歷史以預測 $t+1$。高頻分支以過去 $w=30$ 期的 $\widehat H_t^{KF}$ 建立 supervised windows，並由 M1--M5 各自產生高頻點預測及區間。若 ARIMA 分支的點預測與區間為
\[
    \widehat L_t,\qquad
    [\widehat L_t^{-},\widehat L_t^{+}],
\]
而高頻方法的輸出為
\[
    \widehat H_t,\qquad
    [\widehat H_t^{-},\widehat H_t^{+}],
\]
則最終合併結果定義為
\[
    \widehat Y_t=\widehat L_t+\widehat H_t,
    \qquad
    C_t^{0.95}=
    [\widehat L_t^{-}+\widehat H_t^{-},
     \widehat L_t^{+}+\widehat H_t^{+}].
\]
此合併方式沿用原始研究設定，隱含低頻與高頻不確定性可直接相加之操作性假設。

M1--M4 的 training windows 以 $80\%/20\%$ 隨機切為 model-training 與 validation subsets，training DataLoader 採 batch size $32$ 並打亂 mini-batches；validation 與 testing 不打亂。此程序未直接使用 final-test data，但隨機切割可能使較晚窗口進入 model-training、較早窗口進入 validation，且相鄰滑動窗口可能高度重疊，因此其 validation loss 可能包含 temporal look-ahead bias。本文在原始設定比較中予以保留，以忠實重現原方法；M6 則使用依時間排序的 rolling-origin future-block validation。

\subsection{核心方法概念}
\label{subsec:core-method-concepts}

\subsubsection{DistPred：以樣本表示條件預測分布}

一般 point forecasting model 對每個輸入只輸出單一預測值；DistPred 則令神經網路一次輸出 $K$ 個數值
\[
    \widehat{\boldsymbol H}_t
    =(\widehat H_{t,1},\ldots,\widehat H_{t,K}),
\]
並將其視為條件預測分布的經驗樣本。模型不預先假設誤差必須服從常態分布，而是透過 Continuous Ranked Probability Score（CRPS）同時衡量樣本分布與真值的距離及預測分布本身的離散程度。其經驗形式可寫為
\[
    \widehat{\operatorname{CRPS}}_t
    =\frac{1}{K}\sum_{k=1}^{K}
      |\widehat H_{t,k}-H_t|
    -\frac{1}{2K^2}\sum_{j=1}^{K}\sum_{k=1}^{K}
      |\widehat H_{t,j}-\widehat H_{t,k}|.
\]
第一項懲罰預測樣本偏離真值，第二項避免所有樣本不適當地集中於同一點。模型完成訓練後，可直接由 $K$ 個輸出的經驗分位數取得點預測及 prediction interval；本文取中位數為點預測，並取第 $2.5$ 與 $97.5$ 百分位數為 $95\%$ 區間。M1 與 M2 的差異不在 DistPred loss，而在其 backbone 分別為 ANN 與 Minusformer。

\subsubsection{Quantile Regression：直接估計條件分位數}

Quantile Regression 不估計完整機率密度，而是針對指定的 quantile level $\tau$ 直接學習條件分位數
\[
    Q_{\tau}(H_t\mid X_t).
\]
令預測值為 $\widehat q_{t,\tau}$，其 pinball loss 定義為
\[
    \rho_{\tau}(H_t-\widehat q_{t,\tau})
    =
    \begin{cases}
      \tau(H_t-\widehat q_{t,\tau}),
      & H_t\geq \widehat q_{t,\tau},\\
      (1-\tau)(\widehat q_{t,\tau}-H_t),
      & H_t<\widehat q_{t,\tau}.
    \end{cases}
\]
此非對稱損失使低分位數對高估施加較大懲罰、高分位數對低估施加較大懲罰。本文同時估計 $\tau\in\{0.025,0.5,0.975\}$：$0.5$ 分位數作為點預測，$0.025$ 與 $0.975$ 分位數形成 $95\%$ prediction interval。為避免三條 quantile curves 次序顛倒，M3 與 M4 使用 monotone parameterization 保證
\[
    \widehat q_{t,0.025}
    \leq \widehat q_{t,0.5}
    \leq \widehat q_{t,0.975}.
\]

\subsubsection{Minusformer：以逐層去冗餘強化時間特徵}

Minusformer 為 Transformer 類型之時間序列 backbone，其核心思想是在 encoder layers 中保留輸入流與輸出流，並將各層已解釋的資訊從後續 representation 中扣除。若第 $j$ 層由目前 residual representation $R^{(j)}$ 擷取成分 $O^{(j)}$，則後續輸入可概念化為
\[
    R^{(j+1)}=R^{(j)}-\widehat O^{(j)}.
\]
輸出流亦以逐層相減或累積的方式整合不同層所提取的訊息。這種 ``minus'' 操作的目的，是降低不同 encoder layers 重複學習相同時間特徵的情形，使較深層著重於尚未被前層解釋的 residual structure。最後將各層資訊整合並映射至目標輸出維度。本文不將 Minusformer 本身視為區間方法；它是 M2 與 M4 的共同 backbone，分別搭配 DistPred distribution head 與 Quantile Regression head。

\subsubsection{Sort Bootstrap：保留原始時間索引順序的重抽樣}

傳統 i.i.d. bootstrap 對時間序列觀測值直接有放回抽樣，會打亂時間順序並破壞序列結構。Sort Bootstrap 的做法是先將每一筆觀測與其原始時間索引配對：
\[
    \{(t,H_t):t=1,\ldots,T\}.
\]
每次 bootstrap 對這些成對資料進行 $T$ 次有放回抽樣，得到索引
\[
    (s_1,\ldots,s_T),
\]
再依 $s_j$ 的原始時間索引由小到大排序，並以排序後對應的觀測值建立假樣本。必須注意，此步驟是依時間索引排序，而不是依 $H_t$ 的數值大小排序。此方法使重複抽樣後的一階與二階樣本動差不受排序操作影響，同時嘗試保留原序列由早至晚的結構。本文對每個 Sort Bootstrap 假樣本重新訓練一個 ANN；重複 $100$ 次後，以預測平均作為點估計，並以預測分布的第 $2.5$ 與 $97.5$ 百分位數建立區間。Sort Bootstrap 與 moving-block bootstrap 不同：前者抽取個別時間點後再依索引排序，後者則直接抽取連續時間 blocks。

\subsection{M1：DistPred--ANN}
\label{subsec:m1-method}

M1 使用三層全連接 ANN 作為高頻 predictor，hidden dimensions 為 $(128,256,128)$，各 hidden layer 使用 ReLU 與 dropout $0.2$，輸出 $K=100$ 個 ensemble samples。訓練時將輸出排序，並以 DistPred 的 CRPS--PWM loss 最小化預測分布與真值之差異。learning rate 設為 $0.001$，最多訓練 $50$ epochs，validation patience 為 $5$。測試時以輸出樣本之中位數作為高頻點預測，並以第 $2.5$ 與第 $97.5$ 百分位數形成高頻 $95\%$ 區間。

\subsection{M2：DistPred--Minusformer}
\label{subsec:m2-method}

M2 保留 M1 的 DistPred 分布輸出與 CRPS--PWM loss，但以 Minusformer 取代 ANN backbone。模型使用 label length $15$、$d_{\mathrm{model}}=128$、feed-forward dimension $d_{ff}=128$、$2$ 個 encoder layers、dropout $0.1$ 與 gate parameter $1$，輸出同樣為 $100$ 個 ensemble samples。learning rate、最大 epochs、patience 及分位數取法均與 M1 相同。

\subsection{M3：Quantile--ANN}
\label{subsec:m3-method}

M3 使用與 M1 相同的 ANN hidden dimensions $(128,256,128)$、ReLU、dropout $0.2$ 與 lag window $30$，但輸出改為
\[
    \tau\in\{0.025,0.5,0.975\}
\]
三個條件分位數，並以 pinball loss 訓練。為避免 quantile crossing，模型以 monotone parameterization 保證下界、中位數與上界依序排列。第 $0.5$ 分位數作為點預測，第 $0.025$ 與 $0.975$ 分位數作為高頻區間。learning rate 為 $0.001$，最多 $50$ epochs，patience 為 $5$。

\subsection{M4：Quantile--Minusformer}
\label{subsec:m4-method}

M4 將 M3 的 quantile regression head 接於 Minusformer backbone，並沿用 M2 的 $d_{\mathrm{model}}=128$、$d_{ff}=128$、$2$ 個 encoder layers、label length $15$、dropout $0.1$ 與 gate parameter $1$。模型輸出三個具單調限制的條件分位數，使用 pinball loss、learning rate $0.001$、最多 $50$ epochs 及 patience $5$。

\subsection{M5：Sort Bootstrap--ANN}
\label{subsec:m5-method}

M5 依學長論文所提出之 Sort Bootstrap 執行。首先將 training high-frequency observations 與原始時間索引配對；第 $b$ 次 bootstrap 對
\[
    \{(t,\widehat H_t^{KF})\}_{t=1}^{T}
\]
進行等長度有放回抽樣，再依抽得的原始時間索引排序，而非依觀測值大小排序。排序後的假樣本用 lag window $30$ 建立 ANN training windows。此程序重複 $B=100$ 次，每次皆重新訓練一個 hidden dimensions 為 $(128,256,128)$、ReLU、dropout $0.2$ 的 MSE--ANN；learning rate 為 $0.001$，最多 $50$ epochs，patience 為 $5$。第 $t$ 期的 $100$ 個高頻預測之平均為點預測，第 $2.5$ 與第 $97.5$ 百分位數為高頻區間。

\subsection{六方法之主要原始參數}
\label{subsec:method-parameters}

\begin{table}[htbp]
\centering
\caption{六種方法的主要原始超參數}
\label{tab:original-hyperparameters}
\resizebox{\textwidth}{!}{%
\begin{tabular}{lcccccc}
\hline
方法 & 高頻模型 & lag & hidden/model dimension & 訓練上限 & ensemble/bootstrap & 區間建立方式 \\
\hline
M1 & ANN & 30 & $(128,256,128)$ & 50 epochs & 100 outputs & DistPred percentiles \\
M2 & Minusformer & 30 & $d_{model}=d_{ff}=128$, 2 layers & 50 epochs & 100 outputs & DistPred percentiles \\
M3 & ANN & 30 & $(128,256,128)$ & 50 epochs & --- & quantiles $(.025,.5,.975)$ \\
M4 & Minusformer & 30 & $d_{model}=d_{ff}=128$, 2 layers & 50 epochs & --- & quantiles $(.025,.5,.975)$ \\
M5 & ANN & 30 & $(128,256,128)$ & 50 epochs & 100 resamples & Sort Bootstrap percentiles \\
M6 & ANN & 15 & $(32,16)$ & $\{125,250,500\}$ & 30 hybrid models & OOT adaptive residual interval \\
\hline
\end{tabular}%
}
\end{table}

M6 另使用 $K=4$ 個 rolling-origin future blocks、initial training fraction $0.5$、residual-window candidates $\{100,200,300,400,\mathrm{all}\}$、$101$ 個 shortest-interval $\beta$ candidates、ARIMA BIC search $0\leq p,q\leq4$，以及 Kalman process/measurement variances $(0.5,10)$。

\subsection{評估指標與比較原則}
\label{subsec:evaluation-metrics}

點估計以 MSE、RMSE 與 MAE 評估；區間估計以 empirical coverage、mean width、median width、lower/upper miss rates 與 Gneiting--Raftery interval score 評估。除 coverage 以越接近名目 $95\%$ 為佳外，其餘誤差、寬度與 interval score 越小越佳；然而區間寬度不得脫離 coverage 單獨解讀。每個方法另記錄 end-to-end elapsed time，包含該方法所需的 Kalman/ARIMA、model selection、training、calibration 與 prediction。M1--M5 共用的 Kalman/ARIMA 前處理成本在方法別時間中各自計入，因此方法別時間代表該 pipeline 獨立執行之成本。

\section{研究結果}
\label{sec:results}

\subsection{M1M9--GARCH 模擬資料比較}
\label{subsec:simulation-results}

\subsubsection{資料生成與實驗設定}

模擬資料由低頻 M1 與高頻 M9 相加而成。低頻條件平均為 AR(1)：
\[
    L_t=0.6L_{t-1}+u_t,
    \qquad
    u_t=\sqrt{h_t}z_t,
\]
其中 $z_t\sim N(0,1)$，且
\[
    h_t=0.5+0.15u_{t-1}^{2}+0.80h_{t-1}.
\]
高頻成分遵循 M9 非線性條件平均，innovation standard deviation 固定為 $1$：
\[
    H_t=0.8H_{t-1}
    -\frac{0.8H_{t-1}}{1+\exp(-10H_{t-1})}
    +\varepsilon_t,
    \qquad \varepsilon_t\sim N(0,1).
\]
最終觀測為 $Y_t=L_t+H_t+\eta_t$，$\eta_t\sim N(0,0.15^2)$。每次 Monte Carlo run 重新生成一條序列，training size 為 $650$、test horizon 為 $200$，data seeds 為 2020--2039，共 $20$ 次 paired runs；六種方法在同一 seed 下使用完全相同的觀測路徑。

\subsubsection{原始超參數比較}

表~\ref{tab:simulation-original} 報告各方法保留原始超參數時，$20$ 次 Monte Carlo 的平均表現。M5 的 RMSE 與 MAE 最小，但 coverage 僅 $64.85\%$，顯示其區間過窄；M6 的 coverage 最高，為 $92.20\%$，且點估計誤差僅次於 M5。M3 的 interval score 最低，為 $17.313$，其每 seed 平均時間約 $36.7$ 秒；M6 平均需 $231.3$ 秒，約為 M3 的 $6.3$ 倍。

\begin{table}[htbp]
\centering
\caption{M1M9--GARCH 模擬資料：原始超參數之 20 次 Monte Carlo 平均結果}
\label{tab:simulation-original}
\resizebox{\textwidth}{!}{%
\begin{tabular}{lrrrrrrrrr}
\hline
方法 & MSE & RMSE & MAE & Coverage(\%) & Mean width & Median width & Interval score & 秒/seed & 20-seed 分鐘 \\
\hline
M1 & 11.688 & 3.366 & 2.637 & 90.35 & 11.769 & 11.250 & 17.625 & 36.8 & 12.28 \\
M2 & 12.663 & 3.500 & 2.741 & 89.35 & 12.091 & 11.337 & 18.987 & 37.5 & 12.51 \\
M3 & 11.674 & 3.366 & 2.636 & 91.68 & 12.327 & 11.753 & \textbf{17.313} & \textbf{36.7} & \textbf{12.22} \\
M4 & 13.182 & 3.564 & 2.795 & 66.23 & 7.987 & 7.336 & 35.410 & 38.0 & 12.67 \\
M5 & \textbf{10.516} & \textbf{3.193} & \textbf{2.497} & 64.85 & \textbf{5.783} & \textbf{5.602} & 30.793 & 46.9 & 15.62 \\
M6 & 11.394 & 3.327 & 2.604 & \textbf{92.20} & 12.352 & 12.184 & 17.585 & 231.3 & 77.09 \\
\hline
\end{tabular}%
}
\end{table}

原始設定下各方法的 lower/upper miss rates 依序為：M1 $(4.83\%,4.83\%)$、M2 $(5.60\%,5.05\%)$、M3 $(3.65\%,4.68\%)$、M4 $(15.45\%,18.33\%)$、M5 $(17.60\%,17.55\%)$、M6 $(3.60\%,4.20\%)$。M6 的上下漏失相對平衡，但平均 coverage 仍低於名目 $95\%$，故結果應解讀為「六種方法中最接近名目水準」，而非已達成精確 $95\%$ coverage。

\subsubsection{控制超參數比較}

為降低 predictor capacity 與計算預算差異造成的混淆，另將可共同控制之設定統一為：lag $30$、ANN hidden dimensions $(128,256,128)$、ReLU、learning rate $0.001$、訓練上限 $50$、Kalman variances $(0.5,10)$、ARIMA 搜尋上限參數設為 $5$，以及 M5/M6 ensemble count $100$。惟兩套原始 ARIMA 函數對搜尋上限的端點定義不同：M1--M5 使用 \texttt{range(5)}，實際候選為 $0\leq p,q\leq4$；M6 使用包含端點的搜尋，候選為 $0\leq p,q\leq5$。因此此處屬於「盡量控制」而非所有實作細節完全相同。M6 因 sequential OOT prediction 必須維持 batch size $1$，moving-block bootstrap、rolling-origin cross-fitting 與 adaptive calibration 亦屬方法定義而予以保留。此控制實驗仍保留 M1--M4 原程式的 random training/validation split，故不應解讀為已消除其 validation temporal look-ahead bias。

\begin{table}[htbp]
\centering
\caption{M1M9--GARCH 模擬資料：控制超參數之 20 次 Monte Carlo 平均結果}
\label{tab:simulation-controlled}
\resizebox{\textwidth}{!}{%
\begin{tabular}{lrrrrrrrrr}
\hline
方法 & MSE & RMSE & MAE & Coverage(\%) & Mean width & Median width & Interval score & 秒/seed & 20-seed 分鐘 \\
\hline
M1 & 11.616 & 3.358 & 2.637 & 91.40 & 12.205 & 11.697 & \textbf{17.315} & 34.5 & 11.51 \\
M2 & 12.448 & 3.477 & 2.710 & 89.25 & 11.954 & 11.184 & 19.241 & 35.3 & 11.76 \\
M3 & 11.825 & 3.384 & 2.649 & 90.88 & 11.861 & 11.393 & 17.633 & \textbf{34.4} & \textbf{11.48} \\
M4 & 12.742 & 3.517 & 2.758 & 64.28 & 7.170 & 6.716 & 36.819 & 35.5 & 11.84 \\
M5 & \textbf{10.481} & \textbf{3.189} & \textbf{2.493} & 68.30 & \textbf{6.152} & \textbf{5.971} & 28.506 & 37.5 & 12.51 \\
M6 & 12.260 & 3.449 & 2.696 & \textbf{92.65} & 12.901 & 12.777 & 18.098 & 405.4 & 135.13 \\
\hline
\end{tabular}%
}
\end{table}

控制設定後，M1 的 interval score 最低且點估計誤差僅次於 M5，形成較佳的速度與區間品質折衷；M6 coverage 仍最高，但每 seed 平均時間約為 M1 的 $11.7$ 倍，且 RMSE 與 interval score 均未優於 M1。因此 M6 的主要優勢應定位為 training-only rolling-origin calibration 與較高 coverage，而非計算效率或所有數值指標皆全面較佳。

\subsection{0050 實證資料比較}
\label{subsec:0050-results}

0050 實驗使用 Yahoo Finance 之 0050.TW adjusted close。training period 為 2020-01-01 至 2024-12-31，共 $1215$ 個交易日；test period 為 2025-01-01 至 2026-07-31，共 $382$ 個交易日。六種方法皆在 log adjusted-price scale 建模，最終合併點預測與區間 endpoints 經 exponential transformation 返回價格尺度後評估。此處為單一真實歷史路徑及固定 model seed 2026，故不是不同資料生成路徑的 Monte Carlo 實驗。

\begin{table}[htbp]
\centering
\caption{0050.TW 原始超參數之點預測與區間預測比較}
\label{tab:0050-original}
\resizebox{\textwidth}{!}{%
\begin{tabular}{lrrrrrrrrr}
\hline
方法 & MSE & RMSE & MAE & Coverage(\%) & Mean width & Median width & Interval score & Lower/Upper miss(\%) & 秒 \\
\hline
M1 & 2.177 & 1.476 & 0.950 & 90.58 & 3.986 & 2.945 & 7.543 & 5.50/3.93 & 106.4 \\
M2 & 3.675 & 1.917 & 1.277 & \textbf{91.36} & 5.724 & 4.472 & 8.329 & 6.02/2.62 & 106.7 \\
M3 & 2.170 & 1.473 & 0.950 & 81.41 & 6.513 & 5.812 & 15.134 & 9.42/9.16 & 106.2 \\
M4 & 3.396 & 1.843 & 1.206 & 69.11 & 4.762 & 3.157 & 19.842 & 18.32/12.57 & 108.3 \\
M5 & 4.810 & 2.193 & 1.507 & 74.87 & 3.823 & 3.119 & 16.359 & 5.50/19.63 & 111.5 \\
M6 & \textbf{1.787} & \textbf{1.337} & \textbf{0.881} & 91.10 & \textbf{3.712} & 3.300 & \textbf{7.350} & \textbf{4.71/4.19} & 348.3 \\
\hline
\end{tabular}%
}
\end{table}

在 0050 單一路徑實驗中，M6 同時取得最低 MSE、RMSE、MAE、mean width 與 interval score，且 lower/upper miss rates 相對平衡。M2 的 coverage 最高，但僅比 M6 高約 $0.26$ 個百分點，平均區間則寬約 $54\%$。M6 的執行時間約為 M1 的 $3.3$ 倍。所有方法的 coverage 仍低於名目 $95\%$，因此本結果支持 M6 在此測試期間具有較佳的點預測與區間寬度--漏失折衷，但不構成對其他市場期間或資產之普遍性保證。

\subsection{太陽黑子實證資料比較}
\label{subsec:sunspot-results}

太陽黑子實驗使用月總太陽黑子數（monthly total sunspot number），資料期間為 1900 年 1 月至 2023 年 1 月，共 $1477$ 筆觀測值。資料依時間順序以前 $80\%$ 作為 training set、後 $20\%$ 作為 test set；training period 為 1900 年 1 月至 1998 年 5 月，共 $1181$ 筆，test period 為 1998 年 6 月至 2023 年 1 月，共 $296$ 筆。本節保留六種方法各自的原始模型結構與超參數，直接在原始太陽黑子數值尺度訓練及評估，並固定 model seed 為 2026。由於此處僅使用一條真實歷史路徑及一個固定 seed，因此不是 Monte Carlo 實驗。

\begin{table}[htbp]
\centering
\caption{月總太陽黑子數：原始超參數之點預測與區間預測比較}
\label{tab:sunspot-original}
\resizebox{\textwidth}{!}{%
\begin{tabular}{lrrrrrrrrr}
\hline
方法 & MSE & RMSE & MAE & Coverage(\%) & Mean width & Median width & Interval score & Lower/Upper miss(\%) & 秒 \\
\hline
M1 & 395.486 & 19.887 & 13.672 & 90.20 & 63.794 & 60.566 & 107.028 & 5.41/4.39 & \textbf{109.3} \\
M2 & 397.530 & 19.938 & 13.889 & 92.23 & 64.908 & 58.242 & 103.185 & 3.04/4.73 & 110.1 \\
M3 & 395.454 & 19.886 & 13.745 & 92.57 & 68.512 & 58.151 & \textbf{94.361} & 3.38/4.05 & 108.0 \\
M4 & 414.711 & 20.364 & 13.898 & 63.51 & \textbf{33.093} & \textbf{28.070} & 258.780 & 19.93/16.55 & 111.2 \\
M5 & 384.020 & 19.596 & \textbf{13.356} & 69.59 & 33.614 & 32.134 & 177.780 & 15.20/15.20 & 125.5 \\
M6 & \textbf{383.407} & \textbf{19.581} & 13.915 & \textbf{97.30} & 97.859 & 99.481 & 107.378 & \textbf{1.69/1.01} & 450.6 \\
\hline
\end{tabular}%
}
\end{table}

在原始尺度與原始超參數設定下，M6 取得最低 MSE 與 RMSE，coverage 亦為六種方法中最高，但其 $97.30\%$ coverage 高於名目 $95\%$，且平均區間寬度為六種方法中最大。M3 的 interval score 最低，coverage 為 $92.57\%$，顯示其在本測試期間具有較佳的區間寬度--漏失折衷。M5 的 MAE 最低，惟 coverage 僅為 $69.59\%$；M4 雖有最窄的平均與中位區間寬度，但 coverage 僅為 $63.51\%$，故不能將其窄區間解讀為較佳的區間預測。M6 執行時間約為 M1 的 $4.1$ 倍，其優勢主要在較高 coverage 與較低平方誤差，而代價為較寬的區間及較高計算成本。
